\section{Персептрон и линейные нейронные сети}
\label{sec:perceptron}

\noindent\hspace{1.25cm}\emph{Раздел используется в ЛР~№1 и~№2.}

\subsection{Персептрон Розенблатта}

\emph{Персептрон}, предложенный Ф.~Розенблаттом в 1958~г.~\cite{rosenblatt1958},~--- нейрон
(\ref{eq:neuron}) с пороговой функцией активации. Он относит вход $\mathbf{x}$ к классу~1,
если $\mathbf{w}^{\mathsf T}\mathbf{x} + b \ge 0$, и к классу~0 в противном случае. Геометрически
уравнение
\begin{equation}
\label{eq:hyperplane}
\mathbf{w}^{\mathsf T}\mathbf{x} + b = 0
\end{equation}
задаёт гиперплоскость в пространстве входов (на плоскости~--- прямую $w_1x_1 + w_2x_2 + b = 0$). Вектор
весов $\mathbf{w}$ перпендикулярен этой границе и направлен в сторону класса~1, а смещение $b$ сдвигает
границу относительно начала координат: расстояние от начала координат до прямой равно
$|b| / \lVert\mathbf{w}\rVert$.

\emph{Правило обучения персептрона.} Веса инициализируются нулями или малыми случайными числами. Для
очередного примера $(\mathbf{x}, t)$ вычисляется выход $y$ и ошибка $e = t - y \in \{-1, 0, 1\}$, после
чего веса корректируются:
\begin{equation}
\label{eq:perceptron-rule}
\mathbf{w} \leftarrow \mathbf{w} + \eta\, e\, \mathbf{x}, \qquad b \leftarrow b + \eta\, e ,
\end{equation}
где $\eta > 0$~--- скорость обучения; обычно $\eta = 1$ (при нулевых начальных весах $\eta$ лишь
масштабирует веса и не меняет границу). Если пример
классифицирован верно ($e = 0$), веса не меняются; если точка класса~1 отнесена к классу~0 ($e = 1$),
вектор $\mathbf{w}$ поворачивается к ней, в обратном случае~--- от неё. Эпохи повторяются, пока за
эпоху не будет допущено ни одной ошибки.

\emph{Теорема сходимости} (Розенблатт, Новиков, 1962): если классы \emph{линейно разделимы}, т.\,е.
существует гиперплоскость, строго разделяющая их с зазором $\gamma > 0$, а все входы лежат в шаре
радиуса $R$, то правило (\ref{eq:perceptron-rule}) находит разделяющую гиперплоскость не более чем за
$(R/\gamma)^2$ коррекций независимо от порядка предъявления примеров. Найденная граница не единственна
и зависит от начальных весов и порядка примеров; персептрон ищет \emph{какую-нибудь} разделяющую
прямую, а не <<лучшую>>.

\subsection{Линейная разделимость и проблема исключающего ИЛИ}

Логические функции И и ИЛИ от двух переменных линейно разделимы, а исключающее ИЛИ (XOR) нет: точки
$(0, 1)$ и $(1, 0)$ класса~1 лежат на одной диагонали квадрата, точки $(0, 0)$ и $(1, 1)$ класса~0~---
на другой, и никакая прямая их не разделяет (рисунок~\ref{fig:separability}). Это легко доказать:
из условий $b < 0$, $w_1 + b \ge 0$, $w_2 + b \ge 0$ следует $w_1 + w_2 + b \ge -b > 0$, что противоречит
требованию $w_1 + w_2 + b < 0$ для точки $(1, 1)$.

\begin{figure}[!htb]
\centering
\includegraphics[width=0.96\textwidth]{\figdir/separability.png}
\caption{Линейно разделимые функции И, ИЛИ и неразделимая функция XOR}
\label{fig:separability}
\end{figure}

Для неразделимых данных теорема сходимости не действует: правило персептрона бесконечно исправляет
одну ошибку ценой другой, и число ошибок за эпоху не убывает до нуля. Анализ М.~Минского и
С.~Пейперта (1969)~\cite{minsky1969} показал ограниченность однослойного персептрона и надолго
снизил интерес к нейронным сетям. Решение даёт \emph{скрытый слой}: XOR выражается через разделимые
функции, $x_1 \oplus x_2 = (x_1 \lor x_2) \land \lnot(x_1 \land x_2)$, т.\,е. два нейрона скрытого слоя
строят две прямые (рисунок~\ref{fig:separability}, справа), а выходной нейрон комбинирует их ответы.
Обучать такие сети позволил алгоритм обратного распространения ошибки (раздел~\ref{sec:mlp}).

\subsection{Персептрон с несколькими нейронами}

Слой из $m$ персептронов с общими входами строит $m$ гиперплоскостей, и каждый нейрон отвечает за
один двоичный разряд кода класса. Два нейрона делят плоскость двумя прямыми на четыре области и
различают до $2^2 = 4$ классов, закодированных парами $(y_1, y_2) \in \{00, 01, 10, 11\}$; $m$ нейронов
различают до $2^m$ классов. Правило (\ref{eq:perceptron-rule}) применяется к каждому нейрону независимо
со своей целью~--- соответствующим разрядом кода. Условие применимости: каждый разряд кода должен быть
линейно разделимой функцией входа. Поэтому кодирование классов выбирают так, чтобы соседние на
плоскости группы различались одним разрядом.

\subsection{Линейная нейронная сеть ADALINE}

Линейная сеть (ADALINE, Б.~Уидроу и М.~Хофф, 1960~\cite{widrow1960}) отличается от персептрона
линейной функцией активации: $y = \mathbf{w}^{\mathsf T}\mathbf{x} + b$. Выход непрерывен, поэтому
ошибка $e = t - y$ также непрерывна, и можно минимизировать среднеквадратичную ошибку
(\ref{eq:mse}), которая для линейной модели является квадратичной функцией весов с единственным
минимумом.

\emph{Итерационное обучение} выполняется правилом Уидроу~--- Хоффа (LMS, least mean squares), которое представляет собой стохастический градиентный спуск по мгновенной
квадратичной ошибке:
\begin{equation}
\label{eq:lms}
\mathbf{w} \leftarrow \mathbf{w} + 2\alpha\, e\, \mathbf{x}, \qquad b \leftarrow b + 2\alpha\, e .
\end{equation}
Процесс устойчив при $0 < \alpha < 1/\lambda_{\max}$, где $\lambda_{\max}$~--- наибольшее собственное
число корреляционной матрицы входов $\mathbf{R} = \mathrm{M}[\mathbf{x}\mathbf{x}^{\mathsf T}]$ (в NumPy~---
\code{np.linalg.eigvalsh(P @ P.T / N).max()}).

\emph{Аналитическое решение.} Объединим входные векторы в столбцы матрицы $P$ размера $n \times N$
(добавив строку единиц для смещения), а цели~--- в строку $T$ длины $N$. Минимум суммы квадратов
ошибок $\lVert WP - T\rVert^2$ даётся нормальными уравнениями метода наименьших квадратов (МНК):
\begin{equation}
\label{eq:pinv}
W = T P^{\mathsf T}\bigl(PP^{\mathsf T}\bigr)^{-1} = T P^{+},
\end{equation}
где $P^{+}$~--- псевдообратная матрица Мура~--- Пенроуза. Если $PP^{\mathsf T}$ вырождена (например,
строки $P$ линейно зависимы), $P^{+}$ вычисляется через сингулярное разложение и даёт решение с
наименьшей нормой весов. В NumPy решение даёт \code{W = T @ np.linalg.pinv(P)} или численно более
устойчивая функция \code{np.linalg.lstsq}, решающая задачу МНК без явного обращения матрицы
$PP^{\mathsf T}$.

\subsection{Прогнозирование временного ряда линейной сетью}

Пусть сигнал $x(t)$ дискретизирован с шагом $\Delta t$: $x_i = x(t_i)$, $t_i = t_0 + (i - 1)\Delta t$,
$i = 1, \ldots, N$. Требуется предсказать
следующее значение $x_i$ по $k$ предыдущим $x_{i-1}, \ldots, x_{i-k}$. Для этого на вход линейной сети
подаётся выход \emph{линии задержки} (tapped delay line), и сеть реализует авторегрессионную модель
порядка~$k$:
\begin{equation}
\label{eq:ar}
y_i = w_1 x_{i-1} + w_2 x_{i-2} + \ldots + w_k x_{i-k} + b .
\end{equation}
Матрица входов строится методом скользящего окна: $i$-й столбец содержит $k$ значений сигнала,
предшествующих $x_i$, а целевой вектор~--- сам сигнал:
\begin{equation}
\label{eq:delay-matrix}
P = \begin{pmatrix}
0 & x_1 & x_2 & \cdots & x_{N-1}\\
0 & 0 & x_1 & \cdots & x_{N-2}\\
\vdots & & \ddots & & \vdots\\
0 & 0 & \cdots & \cdots & x_{N-k}
\end{pmatrix}, \qquad T = (x_1, x_2, \ldots, x_N).
\end{equation}
Нули в левой части соответствуют нулевым начальным условиям линии задержки (до начала наблюдения
сигнал считается равным нулю). Веса находятся по формуле (\ref{eq:pinv}).

Для гармонического сигнала $x(t) = \sin(\omega t)$ существует точное рекуррентное соотношение,
вытекающее из формулы суммы синусов:
\begin{equation}
\label{eq:harmonic}
x_i = 2\cos(\omega\Delta t)\, x_{i-1} - x_{i-2}.
\end{equation}
Следовательно, линейная сеть с $k = 2$ уже способна прогнозировать синусоиду без ошибки, а при $k > 2$
МНК находит те же два ненулевых веса. Ошибка возникает только в начале ряда, где линия задержки ещё
заполнена нулями: при $t_0 = 0$ отсчёт $x_2 = \sin(\omega\Delta t)$ не может быть предсказан по
предыстории $(x_1, 0, \ldots, 0) = (0, \ldots, 0)$, и среднеквадратичная ошибка примерно равна $\sin^2(\omega\Delta t)/N$. Она не
зависит от $k$ и уменьшается с шагом дискретизации. Для реальных сигналов с шумом ситуация иная:
увеличение $k$ позволяет усреднять шум по большему числу отсчётов, и ошибка убывает до насыщения
(см. пример выполнения ЛР~№2).

Выбор шага дискретизации ограничен \emph{теоремой Котельникова} (Найквиста~--- Шеннона): сигнал с
наибольшей частотой $f_{\max} = \omega/(2\pi)$ однозначно восстанавливается по отсчётам, если
$\Delta t < 1/(2 f_{\max}) = \pi/\omega$, т.\,е. на период приходится больше двух отсчётов. На практике
для прогнозирования берут 10--20 отсчётов на период: при слишком крупном шаге соседние отсчёты слабо
связаны, модель становится плохо обусловленной и чувствительной к шуму.

\subsection{Реализация на Python}
\label{sec:py-perceptron}

Обучение персептрона~--- два вложенных цикла: по эпохам и по примерам (правило (\ref{eq:perceptron-rule}),
$\eta = 1$). Для многоклассового случая \code{w} становится матрицей, а \code{e}~--- вектором ошибок
нейронов.
\begin{lstlisting}
w, b = np.zeros(2), 0.0
for epoch in range(100):
    errors = 0
    for x, t in zip(X, T):                   # X: (N, 2), T: (N,) из {0, 1}
        e = t - int(w @ x + b >= 0)          # ошибка: -1, 0 или 1
        w, b = w + e * x, b + e
        errors += e != 0
    if errors == 0:                          # эпоха без ошибок: классы разделены
        break
\end{lstlisting}
Для сравнения служит \code{sklearn.linear\_model.Perceptron(eta0=1.0)}: он тоже находит разделяющую
прямую, но из-за перемешивания примеров~--- другую, что иллюстрирует неединственность решения.

Матрица входов линии задержки (\ref{eq:delay-matrix}) строится без циклов функцией
\code{sliding\_window\_view}, веса находятся одним вызовом МНК:
\begin{lstlisting}
from numpy.lib.stride_tricks import sliding_window_view
xp = np.concatenate([np.zeros(k), x])            # нулевые начальные условия
P = sliding_window_view(xp[:-1], k)[:, ::-1].T   # k x N: x[i-1], ..., x[i-k]
Pb = np.vstack([P, np.ones(len(x))])             # строка единиц для смещения
W = np.linalg.lstsq(Pb.T, x, rcond=None)[0]      # то же, что x @ np.linalg.pinv(Pb)
y = W @ Pb                                       # прогноз сети
\end{lstlisting}
Готовые аналоги~--- \code{sklearn.linear\_model.LinearRegression} (МНК) и \code{SGDRegressor}
(итерационное обучение, близкое к правилу LMS).
